\chapter{Method}
\label{chap:analysis-and-design}

This chapter establishes the four scheduler variants and the experimental
protocol used to compare them on the research question: under a fixed
per-generation task budget, which scheduler produces the
highest-performing agent at the lowest cost. All four variants share the
same language model (minimax-m2.5), benchmark, number of generations,
initial archive, seed, and per-generation task budget; differences reside
only in candidate selection and promotion policy, isolating the scheduler
as the sole independent variable. Section~\ref{sec:schedulers} defines
each scheduler and its purpose. Section~\ref{sec:protocol} specifies the
experimental protocol. Section~\ref{sec:metrics} lists the collected
metrics and their relevance to the research question.
Section~\ref{sec:analysis} sets the analysis plan.
Section~\ref{sec:threats} discusses threats to validity.

\section{Subjects Under Comparison}
\label{sec:schedulers}

\paragraph{Baseline.}
Baseline is the base configuration of the system and is used as the
reference scheduler. It evaluates each child through a fixed three-stage
pipeline of cumulative task counts, killing candidates between stages
on a fixed threshold at stage~0 and an archive-relative threshold at
stage~1. The promotion mechanism of this scheduler is based on three
stages: small (3 tasks), medium (12 tasks), and full (35 tasks), for a
total of 50 tasks per fully-evaluated child. To pass from the small to
the medium stage, the child must resolve at least 40\% of the
small-stage tasks. To pass from the medium to the full stage, the
child's cumulative score must meet or exceed the full-evaluation
threshold (the lowest score in the current archive). Once at the full
stage the child is evaluated on the 35 remaining tasks to obtain its
final score.

Shared run controls used by all schedulers:
\texttt{--selfimprove\_size} controls the number of children per
generation; \texttt{--generation\_task\_budget\_total} is the maximum
number of tasks per generation, ensuring a fair budget across schedulers;
\texttt{--selfimprove\_workers} sets the number of evaluations run in
parallel; \texttt{--num\_benchmark\_evals} controls the number of times
each child re-runs the same task subset; \texttt{--benchmark} selects the
benchmark suite (SWE-bench, SWE-bench mini, or polyglot). Other run
controls: seed, max generations, archive update policy, shallow
evaluation, skip-final-evaluation, and evaluation noise.

\paragraph{Hyperband.}
Hyperband is the synchronous Successive Halving implementation for DGM,
where each candidate receives a number of tasks directly proportional to
how well it ranks at each rung. The evaluation is organised into rungs;
each rung evaluates the surviving candidates on its assigned task budget
and keeps the top $1/\eta$ fraction, where $\eta$ is the reduction
factor. $\eta$ controls two things at each rung: the survival fraction
($1/\eta$) and the per-candidate budget growth (approximately
$\eta\times$ larger task budget at the next rung). Per rung: evaluate all
surviving candidates on the incremental tasks, score them by the
cumulative resolution rate across all tasks seen so far, wait for all
evaluations to finish, and promote $\lceil n_r / \eta \rceil$ children
with a positive score. For example, with $n = 15$ and $\eta = 5$:
rung~0 $\rightarrow$ 3 promoted; rung~1 $\rightarrow$ 1 promoted. In the
last rung every surviving candidate is kept.

Hyperband-specific parameters: \texttt{--hyperband\_eta} (the reduction
factor); \texttt{--hyperband\_budgets} (the cumulative task counts at
each rung; default \texttt{2,10,50});
\texttt{--hyperband\_initial\_children} (number of children spawned at
rung~0). This scheduler replaces the baseline's fixed three-stage rule
with a configurable bracket, spending the same task budget on more
candidates and killing the least-promising first.

\paragraph{ASHA.}
ASHA is the asynchronous variant of Hyperband. It uses the same rung
structure and task budgets as Hyperband but never waits for an entire
rung to finish before promoting. The promotion mechanism of ASHA is
conceptually the same as Hyperband but uses a fixed quota per rung
computed at startup: $\max(1, \lceil n / \eta^{r+1} \rceil)$.
Evaluations stream in asynchronously. On each completion at
rung~$r$, ASHA re-ranks all rung-$r$ completions seen so far; if the
candidate has a positive score and currently sits in the top-quota set,
it is promoted to rung~$r+1$ immediately. Otherwise it is dropped and
ASHA continues processing further completions. As soon as a worker slot
frees, the next pending candidate is dispatched, minimising idle time
relative to synchronous Hyperband.

ASHA reuses Hyperband's parameters
(\texttt{--hyperband\bk\_eta}, \texttt{--hyperband\bk\_budgets},
\texttt{--hyperband\bk\_initial\bk\_children}), extending Hyperband
by removing the per-rung synchronisation barrier.

\paragraph{Genetic Algorithm (GA).}
The Genetic Algorithm (GA) scheduler~\cite{Holland:75} is a
blind-evolution control: it preserves DGM's three-stage Baseline
evaluation pipeline but strips error-log context from child
generation. It is used to test whether
DGM's gains come from error-log-guided self-diagnosis or from random
mutations followed by selection. Promotion works exactly the same as
Baseline; the only difference is in child generation, which uses a
high-temperature LLM mutation with no error-log context.

GA-run-specific parameters: \texttt{--ga\_mutation\_temperature}
controls the LLM sampling temperature for blind mutations;
\texttt{--ga\_tournament\_k} controls the tournament
size~\cite{Miller:95} at parent selection (activated by setting
\texttt{--choose\_selfimproves\_method tournament} in
\texttt{DGM\_outer.py}).

\section{Experimental Protocol}
\label{sec:protocol}

The single independent variable across runs is the scheduler choice
(\texttt{--scheduler}), taking one of four values: \texttt{baseline},
\texttt{hyperband}, \texttt{asha}, \texttt{ga}. Four runs were
conducted, one per scheduler.

All other variables are held constant:
\begin{itemize}
  \item \textbf{Language model}: minimax-m2.5 (via OpenRouter).
  \item \textbf{Benchmark}: \texttt{swe\_verified\_mini} (a 50-task
    subset of SWE-bench Verified).
  \item \textbf{Initial archive}: bootstrapped with
    \texttt{test\bk\_swebench.py} on
    \texttt{swe\bk\_verified\bk\_mini} into
    \texttt{initial\bk\_swe\bk\_verified\bk\_mini/}.
  \item \textbf{Random seed}: \texttt{0}.
  \item \textbf{Generations per run}: \texttt{4}
    (\texttt{--max\_generation 4}).
  \item \textbf{Children per generation}: \texttt{2}
    (\texttt{--selfimprove\bk\_size 2}) for Baseline and GA;
    Hyperband and ASHA spawn \texttt{15} instead
    (\texttt{--hyperband\bk\_initial\bk\_children 15}).
  \item \textbf{Per-generation task budget cap}: \texttt{100}
    (\texttt{--generation\bk\_task\bk\_budget\bk\_total 100}).
  \item \textbf{Archive update policy}: \texttt{keep\_all}
    (\texttt{--update\bk\_archive keep\bk\_all}).
  \item \textbf{Benchmark evaluations per child}: \texttt{1}
    (\texttt{--num\_benchmark\_evals 1}).
  \item \textbf{Evaluation noise leeway}: \texttt{0.1}
    (\texttt{--eval\_noise 0.1}).
  \item \textbf{Hardware}: Intel Core i7-12800H (14 cores, 20 threads),
    16~GiB RAM, Fedora Linux 44.
\end{itemize}

A sample run command for Hyperband is given in
Listing~\ref{lst:run-hyperband} (Appendix~\ref{app:reproduction}). The
other three schedulers use identical flags with \texttt{--scheduler}
swapped: see Listings~\ref{lst:run-baseline}, \ref{lst:run-ga},
and~\ref{lst:run-asha} for the Baseline, GA, and ASHA commands. Baseline
drops the three \texttt{--hyperband\_*} flags; GA drops them and adds
\texttt{--ga\_mutation\_temperature 1.0}.

\section{Metrics}
\label{sec:metrics}

The metrics fall into two groups: per-child (recorded in each child's
\texttt{metadata.json}) and per-generation (aggregated and stored in
\texttt{dgm\bk\_metadata.jsonl}). Primary metrics for the research
question are marked \emph{primary}; the rest are \emph{diagnostic}.

\paragraph{Per-child metrics.}
\begin{itemize}
  \item \texttt{accuracy\_score} (\emph{primary}). The rate
    $\mathrm{resolved} / \mathrm{submitted}$ in $[0, 1]$. Empty
    patches count in the denominator (one submitted, zero resolved,
    so they drag the rate down). Tasks with no prediction file at all
    (context-overflow crashes, docker timeouts) drop silently from
    the denominator; see Section~\ref{sec:threats}. Direct measure of
    agent effectiveness on the tasks it actually attempted.
  \item \texttt{is\bk\_compiled} (\emph{diagnostic}). A boolean flag
    (\texttt{is\bk\_compiled\bk\_self\bk\_improve} in
    \texttt{utils/evo\bk\_utils.py}) that is \texttt{True} when the
    child both (a) produced at least one non-empty patch, i.e.\
    $|\mathrm{resolved}| + |\mathrm{unresolved}| > 0$, and (b) has
    \texttt{compile\bk\_failure\bk\_reason} equal to \texttt{None}.
    Tracks wasted budget on children whose patch never compiled or
    was empty.
  \item \texttt{promoted\_from\_rung} (\emph{diagnostic}). Per-child
    integer: the highest rung the child was promoted from
    (\texttt{None} if killed at rung~0). Reveals how aggressively each
    scheduler prunes and how budget is concentrated on survivors.
  \item \texttt{token\bk\_usage\bk\_total} (\emph{primary, cost
    component}). Sum of input + output tokens consumed by the child
    across all LLM calls (mutation generation, evaluation iterations).
    Read directly from
    \texttt{metadata['token\bk\_usage\bk\_total']}. Multiplied by the
    minimax-m2.5 token prices to convert to USD.
  \item \texttt{evaluated\_task\_count} (\emph{diagnostic}). Number of
    distinct tasks the child was actually evaluated on
    (\texttt{total\_submitted\_instances}). Used to determine whether
    the child reached the full budget, and required for inclusion in
    \texttt{best\_child\_score}.
\end{itemize}

\paragraph{Per-generation metrics} (each entry in
\texttt{dgm\_metadata.jsonl}):
\begin{itemize}
  \item \texttt{best\bk\_child\bk\_score} (\emph{primary outcome}).
    $\max(\texttt{accuracy\bk\_score})$ over children whose
    \texttt{evaluated\bk\_task\bk\_count} $\geq$
    \texttt{full\bk\_budget\bk\_size}. Single-number summary of the
    best agent produced in the generation. Tracks DGM's progress over
    generations.
  \item \texttt{children\_fully\_evaluated\_count}
    (\emph{diagnostic}). Number of children whose
    \texttt{evaluated\_task\_count} reached the final rung's task
    budget. Indicates how effective the scheduler was at promoting
    survivors all the way through the pipeline.
  \item \texttt{evaluation\_budget\_tasks\_consumed} (\emph{primary,
    efficiency}). Sum of
    $\mathit{candidates\_at\_rung} \times
    \mathit{tasks\_added\_at\_rung}$
    across all rungs/stages in the generation. The denominator of any
    ``best score per task'' efficiency comparison between schedulers.
  \item \texttt{cost\_per\_gen} (\emph{primary, cost}). Derived (not
    stored):
    $\text{cost} = (\Sigma\ \mathit{input\_tokens}) \times
    \$0.15/\text{M} + (\Sigma\ \mathit{output\_tokens}) \times
    \$1.15/\text{M}$,
    summed across all children in the generation, using
    minimax-m2.5's OpenRouter prices. This is the metric the thesis
    seeks to minimise subject to score gain.
  \item \texttt{generation\_wall\_clock\_seconds}
    (\emph{diagnostic}). Elapsed wall-clock seconds from generation
    start to scheduler return, stored as
    \texttt{generation\_wall\_clock\_seconds}. Captures parallelism
    gains that are invisible to cost and task counts, critical for
    evaluating ASHA's asynchronous claim.
  \item \texttt{archive\_size} and $\Delta$\texttt{archive\_size}
    (\emph{diagnostic}). \texttt{archive\_size} is the cumulative
    count of compiled children in the archive after the generation
    completes; $\Delta$\texttt{archive\_size} is the difference
    between consecutive generations (i.e., the number of new children
    promoted into the archive this generation). Tracks population
    diversity over generations.
\end{itemize}

\section{Analysis Plan}
\label{sec:analysis}

Each scheduler was run once (single seed), so no between-run
aggregation is possible. All analysis is within-run and descriptive,
without significance tests or confidence intervals.

Within each run, the metadata is loaded and the primary metrics
(\texttt{best\bk\_child\bk\_score}, \texttt{cost\bk\_per\bk\_gen},
\texttt{evaluation\bk\_budget\bk\_tasks\bk\_consumed}) are plotted as
a function of generation. Diagnostic metrics
(\texttt{generation\bk\_wall\bk\_clock\bk\_seconds},
\texttt{children\bk\_compiled\bk\_count}, \texttt{archive\bk\_size},
and the \texttt{promoted\bk\_from\bk\_rung} distribution) are plotted
alongside for context. The four schedulers are then compared on the
same axes to
surface differences in trajectory and final state.

The only source of within-run variance is \texttt{--eval\_noise 0.1},
the noise leeway applied to the archive-worst score threshold in
\texttt{update\_archive()}. Score changes below this threshold are not
treated as meaningful improvements.

Aggregated cost, wall-clock, final archive composition, and
best-of-run score per scheduler are reported in
Chapter~\ref{chap:results}, Section~\ref{sec:results-summary}. The best-of-run score is
defined as $\max(\texttt{best\_child\_score})$ across the four
generations of a scheduler's run that is, the best child produced
at any point in the run. Per-generation
trajectories appear in Section~\ref{sec:results-accuracy}. The
cost-vs-score comparison, which is the central result for the research
question, appears in Section~\ref{sec:results-costscore}.

A multi-seed re-run of all four schedulers is required before any
inferential claim can be made; this is deferred to future work.

\section{Threats to Validity}
\label{sec:threats}

Several threats apply to this work.

\paragraph{Internal validity.}
\begin{itemize}
  \item Single-seed-per-scheduler execution: one run per condition, no
    within-scheduler variance estimate, no significance claim possible.
  \item \texttt{accuracy\_score} denominator silently drops tasks for
    which no prediction file was emitted (context-overflow crashes,
    docker timeouts, harness exceptions). Empty patches are counted
    and penalise the rate, but full no-file failures shrink the
    denominator and can inflate the reported rate; minimax-m2.5
    context-overflow crashes on harder tasks expose this bias.
  \item Hyperband and ASHA require score $> 0$ for promotion, killing
    zero-score children regardless of their tie-breaker rank.
  \item \texttt{num\_benchmark\_evals = 1}: each child's score is a
    single-sample estimate of LLM-stochastic performance.
\end{itemize}

\paragraph{External validity.}
\begin{itemize}
  \item Results on the 50-task \texttt{swe\_verified\_mini} subset may
    not generalise to the full SWE-bench Verified (500 tasks), to other
    SWE benchmarks, or to polyglot tasks.
  \item All four runs use minimax-m2.5; scheduler behaviour under
    stronger or weaker LLMs may differ.
  \item Each run is capped at four generations; long-run dynamics are
    not observed.
  \item All four runs start from the same initial archive; biases in
    the seed agent propagate uniformly.
\end{itemize}

\paragraph{Construct validity.}
\begin{itemize}
  \item \texttt{cost\_per\_gen} measures LLM-spend only;
    Docker/infrastructure cost is excluded.
  \item Cost-vs-score framing assumes uniform task value, which
    SWE-bench Verified does not enforce.
  \item Wall-clock comparisons are confounded by background system load
    and OpenRouter latency, particularly relevant for ASHA's
    asynchronous claim.
\end{itemize}
